\begin{derivation}[从\eqnrefs{eq:qcd-qqbar-gg-amplitude-sum-dp68}到\eqnrefs{eq:qcd-qqbar-gg-ward-dp32}]
把两条出射胶子动量记为 $k_1,k_2$，颜色指标为 $a,b$。将三张图共同的 $g_s^2$ 因子提出，并定义约化张量
\begin{equation*}
 \widetilde{\mathcal M}_x^{\mu\nu}
 \equiv g_s^{-2}\mathcal M_x^{\mu\nu},
 \qquad x=t,u,s.
\end{equation*}
对两张夸克交换图，费米子线结构严格写成
\begin{align*}
 \widetilde{\mathcal M}_t^{\mu\nu}
 &=\bar v(p_2)\gamma^\nu T^b
 S(p_1-k_1)\gamma^\mu T^a u(p_1),\\
 \widetilde{\mathcal M}_u^{\mu\nu}
 &=\bar v(p_2)\gamma^\mu T^a
 S(p_1-k_2)\gamma^\nu T^b u(p_1),
\end{align*}
其中 $S(q)$ 表示 Dirac 算符 $S^{-1}(q)=\slashed q-mc$ 的 Feynman 逆核。代数运算先在内部费米子分母非零的开运动学区域进行，所得恒等式再按统一的 $+\ii0$ 处方延拓。

把第一条胶子偏振替换为 $k_1^\mu$。由
\begin{equation*}
 \slashed k_1
 =S^{-1}(p_1)-S^{-1}(p_1-k_1)
\end{equation*}
以及在壳方程 $S^{-1}(p_1)u(p_1)=0$，得到
\begin{align*}
 S(p_1-k_1)\slashed k_1u(p_1)
 &=S(p_1-k_1)
 [S^{-1}(p_1)-S^{-1}(p_1-k_1)]u(p_1)\\
 &=-u(p_1).
\end{align*}
因此 $k_{1\mu}\widetilde{\mathcal M}_t^{\mu\nu}$ 留下一个没有内部传播子的费米子流，颜色顺序为 $T^bT^a$。

对另一种排序，利用四动量守恒 $p_1+p_2=k_1+k_2$，把 $k_1$ 写成靠近反费米子端的两个逆传播子之差，并令逆传播子向左作用到 $\bar v(p_2)$。由 $\bar v(p_2)(\slashed p_2+mc)=0$，得到同一个旋量流，颜色顺序为 $T^aT^b$，符号与上一项相反。两张约化夸克图的纵向收缩之和因而含有
\begin{equation*}
 T^aT^b-T^bT^a=[T^a,T^b]=\ii f^{abc}T^c
\end{equation*}
乘同一个费米子流。阿贝尔理论的生成元是普通数，所以这一对易子严格为零；QCD 的非零对易子要求三胶子图参与同一 Ward 恒等式。

$s$ 道图含夸克颜色流
\begin{equation*}
 J_c^\rho=\bar v(p_2)\gamma^\rho T^c u(p_1)
\end{equation*}
以及由 Yang--Mills 三次作用量得到的三胶子张量 $V_{\rho}{}^{\mu\nu}$。按“全部动量流入”的统一约定收缩 $k_{1\mu}V^{\rho\mu\nu}$，再使用 $k_1^2=k_2^2=0$、$k_2\cdot\varepsilon_2=0$ 与 $q=k_1+k_2$，得到
\begin{equation*}
 k_{1\mu}\widetilde{\mathcal M}_s^{\mu\nu}
 =-k_{1\mu}
 \left(\widetilde{\mathcal M}_t^{\mu\nu}
 +\widetilde{\mathcal M}_u^{\mu\nu}\right).
\end{equation*}
恢复共同因子 $g_s^2$ 后，三图满足
\begin{equation*}
 k_{1\mu}
 \left(\mathcal M_t^{\mu\nu}
 +\mathcal M_u^{\mu\nu}
 +\mathcal M_s^{\mu\nu}\right)=0.
\end{equation*}
交换 $1\leftrightarrow2$ 同样得到第二条 Ward 恒等式。三图抵消依次使用 Dirac 逆传播子的望远镜恒等式、颜色生成元的对易关系和由同一 Yang--Mills 作用量固定的三胶子顶角。
\end{derivation}
